\documentclass[11pt,a4paper]{article}
\usepackage[utf8]{inputenc}
\usepackage[T1]{fontenc}
\usepackage[spanish]{babel}
\usepackage{amsmath, amsthm, amssymb}
\usepackage{xcolor}
 \usepackage{tikz}
% --- Tipografía sans-serif ---
\usepackage{lmodern}
\renewcommand{\familydefault}{\sfdefault}
 
% Gráficos 2D y 3D
\usepackage{tikz}
\usepackage{pgfplots}
\pgfplotsset{compat=1.18, colormap/viridis}
 
% Cuadros personalizados, sombras y entornos numerados
\usepackage[many]{tcolorbox}
\usepackage{multirow}
\usepackage{capt-of}  % \captionof: pies de figura sin flotantes
\tcbuselibrary{skins, breakable, theorems}

\usetikzlibrary{babel, decorations.markings, decorations.pathreplacing, arrows.meta}
\usepackage{etoc} 

% --- Modo exportación (para subir cada tema por separado) ---
% Solo se activa con ./exportar.sh, que compila con \def\exportar{}. En una compilación normal no cambia nada.
\newif\ifexportar
\ifdefined\exportar \exportartrue \else \exportarfalse \fi

\newcommand{\portadaseccion}[1]{
    \clearpage
    \thispagestyle{empty}
    \vspace*{2cm}
    \section{#1} 
    \vspace{1.5cm} 
    \etocsettocstyle{\Large\bfseries Contenido del tema\par\vspace{0.5cm}}{}
    \localtableofcontents
    \clearpage 
}
 
% --- REFERENCIAS Y CITAS ENLAZADAS ---
\usepackage{hyperref}
\usepackage[capitalise, spanish]{cleveref}
\hypersetup{
    colorlinks   = true,
    linkcolor    = blue!75!black,
    citecolor    = violet!60!black,
    urlcolor     = cyan!70!black,
    pdftitle     = {Apuntes de Topología Algebraica},
    pdfborder    = {0 0 0},
    bookmarksnumbered = true,
}
 
% Nombres para que \cref reconozca las cajas personalizadas
\crefname{tcb@cnt@definicion}{definición}{definiciones}
\crefname{tcb@cnt@observacion}{observación}{observaciones}
\crefname{tcb@cnt@nota}{nota}{notas}
\crefname{tcb@cnt@teorema}{teorema}{teoremas}
\crefname{tcb@cnt@proposicion}{proposición}{proposiciones}
\Crefname{tcb@cnt@definicion}{Definición}{Definiciones}
\Crefname{tcb@cnt@observacion}{Observación}{Observaciones}
\Crefname{tcb@cnt@nota}{Nota}{Notas}
\Crefname{tcb@cnt@teorema}{Teorema}{Teoremas}
\Crefname{tcb@cnt@proposicion}{Proposición}{Proposiciones}
 
% --- Fondo y texto ---
\pagecolor{gray!5}
\color{black}
 
% --- Línea de degradado ("scan line") bajo cada \section (CORREGIDA) ---
\usepackage{titlesec}
\newcommand{\neonrule}[1]{%
    \par\noindent % Fuerza el origen al margen izquierdo
    \begin{tikzpicture}[remember picture, overlay]
    % 3 Bloques izquierdos
    \fill[#1] (0,0.3) rectangle (0.2,0.45);
    \fill[#1] (0.25,0.3) rectangle (0.45,0.45);
    \fill[#1] (0.5,0.3) rectangle (0.7,0.45);
    
    % Línea central 
    \draw[#1, line width=1pt] (1.2,0.375) -- (\dimexpr\linewidth-0.4cm\relax,0.375);
    
    % Bloque derecho
    \fill[#1] (\dimexpr\linewidth-0.5cm\relax, 0.3) rectangle (\linewidth,0.45);
\end{tikzpicture}%
}
\titleformat{\section}
  {\normalfont\Large\bfseries}{\thesection}{1em}{}
  [\neonrule{violet}] % Color lila para la línea de la sección
\titlespacing*{\section}{0pt}{3.5ex plus 1ex minus .2ex}{2.3ex plus .2ex}
 
% --- Sombra con brillo de color + título en degradado ---
\tcbset{
    futuristicbase/.style={
        enhanced, breakable,
        boxrule=0.5pt,
        arc=2pt,
        fonttitle=\bfseries
    }
}
 
% --- Modo borrador: sin sombras difusas (son lo más lento de compilar) ---
% En local compila en borrador; GitHub Actions define \versionfinal y compila con sombras.
% Para la versión final en local, pon \def\versionfinal{} en la primera línea de main.tex.
\newif\ifborrador
\ifdefined\versionfinal \borradorfalse \else \borradortrue \fi
\ifborrador
  \tcbset{sombra/.style={}}
\else
  \tcbset{sombra/.style={fuzzy shadow={1.2mm}{-1.2mm}{0mm}{0.3mm}{#1}}}
\fi

% --- Cuadros personalizados numerados por sección (Paleta Azul y Lila) ---
 
\newtcbtheorem[number within=section]{teorema}{Teorema}{
    futuristicbase,
    colback=blue!3!white,
    colframe=blue!80!black,
    sombra={blue!50!black,opacity=0.5},
    title style={left color=blue!90!black, right color=blue!60!black}
}{thm}

\newtcbtheorem[number within=section]{definicion}{Definición}{
    futuristicbase,
    colback=violet!4!white,
    colframe=violet!80!black,
    sombra={violet!50!black,opacity=0.5},
    title style={left color=violet!80!black, right color=violet!50!black}
}{def}

\newtcbtheorem[number within=section]{proposicion}{Proposición}{
    futuristicbase,
    colback=blue!2!violet!3!white,
    colframe=blue!50!violet!50!black,
    sombra={blue!40!violet!40!black,opacity=0.5},
    title style={left color=blue!60!violet!60!black, right color=blue!40!violet!40!black}
}{prop}

\newtcbtheorem[number within=section]{observacion}{Observación}{
    futuristicbase,
    colback=purple!5!white,
    colframe=purple!50!black,
    boxrule=0.4pt,
    coltitle=black,
    sombra={purple!30!black,opacity=0.5},
    title style={left color=purple!20!white, right color=purple!5!white}
}{obs}

\newtcbtheorem[number within=section]{nota}{Nota}{
    futuristicbase,
    colback=magenta!5!blue!5!white,
    colframe=magenta!60!blue!60!black,
    boxrule=0.4pt,
    coltitle=white,
    sombra={magenta!40!blue!40!black,opacity=0.5},
    title style={left color=magenta!80!blue!80!black, right color=magenta!60!blue!60!black}
}{not}

% Ejercicios resueltos en clase: \begin{ejercicio}{Título (p. ej. "Tema 5, ej. 3")} enunciado y solución \end{ejercicio}
% Sin contador: el título lleva la numeración de la hoja de ejercicios.
\newtcolorbox{ejercicio}[1]{
    futuristicbase,
    colback=teal!4!white,
    colframe=teal!60!black,
    boxrule=0.4pt,
    coltitle=white,
    sombra={teal!40!black,opacity=0.5},
    title style={left color=teal!70!black, right color=teal!50!black},
    title={Ejercicio: #1},
    before upper={\ifexportar\phantomsection\addcontentsline{toc}{subsubsection}{#1}\fi}
}

% --- Cabecera/pie de página ---
\usepackage{fancyhdr}
\setlength{\headheight}{12.5pt}
\pagestyle{fancy}
\fancyhf{}
\renewcommand{\headrulewidth}{0.4pt}
\renewcommand{\footrulewidth}{0.4pt}
\fancyhead[L]{\small\textcolor{gray}{\leftmark}}
\fancyhead[R]{\small\textcolor{gray}{Apuntes}}
\fancyfoot[C]{\small\textcolor{gray}{\thepage}}
 

% --- Marcas de los resultados demostrados en clase (⊙ en "Contenidos de las clases") ---
% T y C destacadas; P y L discretas.
\newcommand{\marcadestacada}[1]{\texorpdfstring{\hspace{0.4em}\fcolorbox{white}{red!80!black}{\textcolor{white}{\textbf{$\odot$\,#1}}}}{}}
\newcommand{\marcasuave}[1]{\texorpdfstring{\hspace{0.4em}\fcolorbox{white}{white}{\textcolor{black}{$\odot$\,#1}}}{}}
\newcommand{\marcaT}{\marcadestacada{T}}
\newcommand{\marcaC}{\marcadestacada{C}}
\newcommand{\marcaP}{\marcasuave{P}}
\newcommand{\marcaL}{\marcasuave{L}}

% --- Ejercicios propuestos en clase ("Ejercicio:" en "Contenidos de las clases"), demostrados en el texto ---
\newcommand{\marcaE}{\texorpdfstring{\hspace{0.4em}\fcolorbox{white}{orange!85!black}{\textcolor{white}{\textbf{Ej.\ de clase}}}}{}}

% --- Parte de ejercicios: procedencia de cada solución (prioridad: apuntes > clase > examen > propia) ---
\newcommand{\etiqueta}[2]{\fcolorbox{white}{#1}{\textcolor{white}{\small\textbf{#2}}}}
\newcommand{\fuenteclase}[1]{\etiqueta{teal!70!black}{Resuelto en clase (#1)}}
\newcommand{\fuenteexamen}[1]{\etiqueta{blue!70!black}{Solución del profesor: examen #1}}
\newcommand{\fuentepropia}{\etiqueta{gray!70!black}{Solución propia}}
\newcommand{\enexamen}[1]{\etiqueta{red!80!black}{Cayó en examen: #1}}
\newcommand{\sol}{\par\medskip\noindent\textbf{Solución.}\ }
% Solución que ya está en los apuntes: solo se remite a ella
\newcommand{\enapuntes}[1]{\par\medskip\noindent\textbf{Solución.}\ \etiqueta{teal!70!black}{En los apuntes} #1}

% --- Partes exportables ---
% \inicioparte{id}{encabezado}{título}{profundidad del índice}: en modo exportación empieza una parte nueva
% con portadilla y hoja de índice (si la profundidad está vacía, no hay índice) y anota su página inicial
% en \jobname.partes para que exportar.sh corte el PDF. \fronteraparte{id} solo anota dónde acaba lo anterior.
% \indicetema{id} va justo después del \section o \subsection de la parte y guarda su índice local.
\ifexportar
  \newwrite\exportpartes
  \immediate\openout\exportpartes=\jobname.partes
\fi
\newcommand{\inicioparte}[4]{%
  \ifexportar
    \clearpage
    \thispagestyle{empty}%
    \write\exportpartes{#1 \thepage}%
    \vspace*{5cm}
    \begin{center}
      {\Large\bfseries Topología Algebraica\par}
      \vspace{0.6em}
      {\normalsize Grado en Matemáticas \textperiodcentered\ Curso 2026/27\par}
      \vspace{3cm}
      {\Huge\bfseries #2\par}
      \vspace{1em}
      {\LARGE #3\par}
    \end{center}
    \clearpage
    \if\relax\detokenize{#4}\relax\else
      \thispagestyle{empty}%
      {\Large\bfseries Índice\par}
      \vspace{1em}
      {\etocsettocstyle{}{}\etocsetnexttocdepth{#4}\tableofcontents \ref{toc:#1}}
      \clearpage
    \fi
  \fi}
\newcommand{\fronteraparte}[1]{\ifexportar\write\exportpartes{#1 \thepage}\fi}
\newcommand{\indicetema}[1]{\ifexportar\invisiblelocaltableofcontents\label{toc:#1}\fi}
